%=== ABSTRACT ===
\newpage
\chapter*{\centering Abstract}
\markboth{Abstract}{}
\addcontentsline{toc}{chapter}{Abstract}

\begin{spacing}{1.5}
\setlength{\parskip}{0.18in}

Delivery quadrotors change mass and roll and pitch inertia abruptly when they
grasp or release a parcel, and a predictive controller that keeps the unloaded
model plans with the wrong hover thrust and the wrong rotational response. This
dissertation studies how far such a vehicle can adapt to a centred rigid payload
from ordinary flight signals, and how it should decide that a commanded release
has physically happened.

For a payload grasped on the thrust axis, rigid-body algebra makes the inertia
increment the reduced mass times the square of a nominal vertical arm, so the
composite mass is the only payload quantity to be identified online. A moving horizon estimator
identifies it from odometry and the thrust reconstructed from measured rotor
speeds. In its deployed form it keeps only velocity and mass as states and takes
the measured attitude as a known input, so that the less reliable rotational
model has no path to the mass. In a paired closed-loop comparison of $96$
flights against a full-state estimator, this form lowered the loaded mass error
from \SI{0.035}{\kg} to \SI{0.013}{\kg}, removed the collapse of the estimate in
manoeuvre and solved about ten times faster, without changing tracking or release
safety. Replayed on the same flights, an extended Kalman filter with the same
inputs was as accurate: the gain comes from the model structure, not from the
moving horizon. One
health-gated estimate updates a nonlinear model predictive controller
continuously; a separate lifecycle accepts a release only when evidence
confirms that the payload has left --- in the deployed form, a force balance of
the reconstructed thrust against the specific force measured by the inertial
unit, which does not pass through the mass estimate.

The system is evaluated in PX4/Gazebo software-in-the-loop simulation with
acados solvers. The deployed combination of the reduced estimator and the
force-balance rule was validated in $180$ flights from a tagged commit, at two
payloads and two speeds, under nominal and delayed release, uncommanded loss, a
jammed gripper and a \SI{\pm10}{\percent} thrust-map error; four acceptance
criteria inherited from an earlier pre-specification were fixed during the run
and all met. All $72$ nominal releases were confirmed, none early and none during
carry, \SI{2.2}{\second} after physical detachment at the median; the loaded mass
error was \SI{0.013}{\kg}, the estimate never collapsed in manoeuvre, and the
estimator solved in \SI{0.3}{\milli\second}. All $24$ loads lost without a
command were detected and none of $24$ jammed releases was accepted. The thrust
map bounds the system: \SI{10}{\percent} high, every release was missed and
ended in a safe hover; \SI{10}{\percent} low, the \SI{0.15}{\kg} payload read as
absent, the loss watchdog never armed, and with detachment delayed by
\SI{4}{\second} all $12$ releases were accepted about \SI{3}{\second} before the
payload left.

Historical batches with the full-state estimator show how the design was
reached. In a frozen cohort of $120$ flights the mass estimate was within
\SI{0.5}{\percent} in loaded hover but collapsed periodically in manoeuvre at
\SI{4}{\meter\per\second}, and confirming the release on that mass evidence
throughout carry produced $114$ false acceptances in $20$ of $32$ flights;
arming it only after the release command removed them. A campaign of $312$
flights introduced the force-balance rule, which removed them without that
restriction, and showed that estimating the mass gave no consistent tracking
benefit over a residual specific force alone.

The evidence is from simulation only, and no hardware validity is claimed. Both
the estimator and the release detector depend on thrust-map calibration, and
neither was compared against a change-detection baseline.

\textbf{Keywords:} quadrotor; nonlinear model predictive control; moving horizon
estimation; online mass estimation; payload inertia; release confirmation;
software-in-the-loop.
\end{spacing}
\newpage
%=== END OF ABSTRACT ===
